\chapter{2D Frame Element}
\label{ch:frame}

% =============================================================================
\section{Frame Element}
\label{sec:frame_element}
% =============================================================================

A \emph{frame element} (beam-column) combines the axial stiffness of a bar
element with the bending stiffness of an Euler-Bernoulli beam element into a
single structural member. Unlike truss elements, which transmit axial force
only through pinned joints, frame elements transmit \emph{axial force, shear
force, and bending moment} at rigid joints. Frames are the appropriate model
whenever joint rigidity induces bending in the members: portal frames,
multi-storey building frames, and any structure in which loads are not applied
along member axes.

In the \emph{local coordinate system} $(\bar{x},\bar{y})$ aligned with the
element axis, each node carries three DOFs --- axial displacement $\bar{u}$,
transverse displacement $\bar{v}$, and rotation $\theta$ (positive
counter-clockwise). The six-component local DOF vector and its conjugate
nodal force vector are:
\begin{equation}
    \bar{\vect{u}}^e =
    [\bar{u}_1,\,\bar{v}_1,\,\theta_1,\,
     \bar{u}_2,\,\bar{v}_2,\,\theta_2]\transpose, \qquad
    \bar{\vect{f}}^e =
    [\bar{N}_1,\,\bar{V}_1,\,\bar{M}_1,\,
     \bar{N}_2,\,\bar{V}_2,\,\bar{M}_2]\transpose,
    \label{eq:frame_dof}
\end{equation}
where $\bar{N}_i$ is the axial force, $\bar{V}_i$ the shear force, and
$\bar{M}_i$ the bending moment at node~$i$ (see \cref{fig:frame_element}).

\begin{figure}[htbp]
\centering
\begin{tikzpicture}[>=Stealth, x=1cm, y=1cm,
    disp/.style={->, very thick, blue!70!black},
    force/.style={->, ultra thick, red!75!black}]

  %% Element body
  \fill[gray!22] (0,-0.22) rectangle (7,0.22);
  \draw[thick]   (0,-0.22) rectangle (7,0.22);
  \node[above=4pt] at (3.5,0.22) {$E,\;A,\;I$};

  %% Nodes
  \filldraw (0,0) circle (3.5pt);
  \filldraw (7,0) circle (3.5pt);
  \node[below left=2pt]  at (0,0) {\small\bfseries 1};
  \node[below right=2pt] at (7,0) {\small\bfseries 2};

  %% Local axis label
  \draw[->,thin] (-0.3,-0.7) -- ++(1.8,0) node[right] {$\bar{x}$};
  \draw[->,thin] (-0.3,-0.7) -- ++(0,1.2) node[above] {$\bar{y}$};

  %% Displacements at node 1 (blue, above)
  \draw[disp] (-0.8,0.5) -- ++(0.8,0) node[above,midway] {$\bar{u}_1$};
  \draw[disp] (0,0.65)   -- ++(0,0.9) node[right=2pt,midway] {$\bar{v}_1$};
  \draw[disp,->] (0.05,1.72) arc[start angle=200,end angle=70,radius=0.55]
        node[above left=-2pt] {$\theta_1$};

  %% Displacements at node 2 (blue, above)
  \draw[disp] (7.0,0.5) -- ++(0.8,0) node[above,midway] {$\bar{u}_2$};
  \draw[disp] (7.0,0.65) -- ++(0,0.9) node[right=2pt,midway] {$\bar{v}_2$};
  \draw[disp,->] (7.05,1.72) arc[start angle=200,end angle=70,radius=0.55]
        node[above left=-2pt] {$\theta_2$};

  %% Forces at node 1 (red, below)
  \draw[force] (-0.9,-0.55) -- ++(0.9,0) node[below,midway] {$\bar{N}_1$};
  \draw[force] (0,-0.65)    -- ++(0,-0.9) node[right=2pt,midway] {$\bar{V}_1$};
  \draw[force,->] (0.05,-1.57) arc[start angle=160,end angle=290,radius=0.55]
        node[below left=-2pt] {$\bar{M}_1$};

  %% Forces at node 2 (red, below)
  \draw[force] (7.0,-0.55) -- ++(0.9,0) node[below,midway] {$\bar{N}_2$};
  \draw[force] (7,-0.65)   -- ++(0,-0.9) node[right=2pt,midway] {$\bar{V}_2$};
  \draw[force,->] (7.05,-1.57) arc[start angle=160,end angle=290,radius=0.55]
        node[below left=-2pt] {$\bar{M}_2$};

  %% Length dimension
  \draw[<->,thin] (0,-2.40) -- (7,-2.40)
        node[fill=white,inner sep=1pt,midway] {$L$};

\end{tikzpicture}
\caption{Two-node frame element in the local frame $(\bar{x},\bar{y})$.
Displacements $\bar{u}_i, \bar{v}_i$ (blue) and rotations $\theta_i$
(blue, CCW positive) are the kinematic DOFs. Conjugate nodal forces
$\bar{N}_i$ (axial), $\bar{V}_i$ (shear), and moments $\bar{M}_i$ (red)
are the static conjugates.}
\label{fig:frame_element}
\end{figure}

% =============================================================================
\section{Local Stiffness Matrix}
\label{sec:frame_local_k}
% =============================================================================

In the local frame, axial deformation (bar action) and transverse bending
(beam action) are decoupled: an axial displacement $\bar{u}$ produces no
bending forces, and a transverse displacement $\bar{v}$ produces no axial
force. The local stiffness matrix $\bar{\matr{k}}^e$ is therefore
block-diagonal, assembling the bar stiffness (\cref{eq:bar_stiffness})
and the Euler-Bernoulli beam stiffness (\cref{eq:beam_stiffness}) into a
single $6\times6$ matrix:

\begin{equation}
    \bar{\matr{k}}^e =
    {\renewcommand{\arraystretch}{1.6}%
    \begin{bmatrix}
        \dfrac{EA}{L}   & 0 & 0 & -\dfrac{EA}{L} & 0 & 0 \\
        0 & \dfrac{12EI}{L^3} & \dfrac{6EI}{L^2} & 0 & -\dfrac{12EI}{L^3} & \dfrac{6EI}{L^2} \\
        0 & \dfrac{6EI}{L^2} & \dfrac{4EI}{L} & 0 & -\dfrac{6EI}{L^2} & \dfrac{2EI}{L} \\
        -\dfrac{EA}{L} & 0 & 0 & \dfrac{EA}{L} & 0 & 0 \\
        0 & -\dfrac{12EI}{L^3} & -\dfrac{6EI}{L^2} & 0 & \dfrac{12EI}{L^3} & -\dfrac{6EI}{L^2} \\
        0 & \dfrac{6EI}{L^2} & \dfrac{2EI}{L} & 0 & -\dfrac{6EI}{L^2} & \dfrac{4EI}{L}
    \end{bmatrix}}
    \label{eq:frame_local_k}
\end{equation}

\begin{importantbox}
The frame local stiffness is block-diagonal: rows and columns $\{1,4\}$
carry the bar block ($EA/L$); rows and columns $\{2,3,5,6\}$ carry the
beam block ($EI/L^3$ scaling). Axial and bending deformations are
\emph{uncoupled} in the local frame.
\end{importantbox}

\begin{remark}
The unconstrained frame element has a three-dimensional null space:
rigid translation in $\bar{x}$, rigid translation in $\bar{y}$, and rigid
rotation. At least three linearly independent kinematic constraints are
needed to make the assembled stiffness non-singular.
\end{remark}

% =============================================================================
\section{Element Load Vector}
\label{sec:frame_load}
% =============================================================================

When distributed loads act along the element, the consistent nodal load
vector $\bar{\vect{f}}^e$ is formed by superposing the bar and beam
contributions derived in \cref{ch:1d,ch:beam}. For a uniform \emph{axial}
distributed load $p_0$ (positive in $+\bar{x}$) and a uniform
\emph{transverse} distributed load $q_0$ (positive in $+\bar{y}$):
\begin{equation}
    \bar{\vect{f}}^e
    =
    \underbrace{
      \frac{p_0 L}{2}
      \begin{bmatrix}1\\0\\0\\1\\0\\0\end{bmatrix}
    }_{\text{axial}\;(\cref{eq:bar_load})}
    +
    \underbrace{
      \frac{q_0 L}{12}
      \begin{bmatrix}0\\6\\L\\0\\6\\{-L}\end{bmatrix}
    }_{\text{transverse}\;(\cref{eq:beam_load})}
    \label{eq:frame_load}
\end{equation}
The two contributions act on separate DOF subsets: the axial terms go to
DOFs~1 and~4 ($\bar{u}$ displacements); the transverse terms go to DOFs~2,
3, 5, and~6 ($\bar{v}$ displacements and $\theta$ rotations).
In the worked example of \cref{sec:frame_example}, the beam (element~2)
carries a downward transverse load, so $q_0 = -q < 0$ and $p_0 = 0$;
the two columns carry no distributed load.

% =============================================================================
\section{Coordinate Transformation}
\label{sec:frame_transform}
% =============================================================================

Frame members are generally inclined in the global $(x,y)$ plane. The
$6\times6$ transformation matrix $\matr{T}_6$ maps from the global frame to
the local frame. With $c = \cos\alpha$ and $s = \sin\alpha$, where $\alpha$
is the element inclination angle:

\begin{equation}
    \matr{T}_6 =
    \begin{bmatrix}
         c &  s & 0 & 0 & 0 & 0 \\
        -s &  c & 0 & 0 & 0 & 0 \\
         0 &  0 & 1 & 0 & 0 & 0 \\
         0 &  0 & 0 & c & s & 0 \\
         0 &  0 & 0 &-s & c & 0 \\
         0 &  0 & 0 & 0 & 0 & 1
    \end{bmatrix}
    \label{eq:frame_T6}
\end{equation}

$\matr{T}_6$ consists of two identical $3\times3$ rotation blocks (one per
node), each rotating the displacement pair $(u_i,v_i)$ from global to local
while leaving the rotation $\theta_i$ unchanged.  The structure mirrors the
truss transformation (\cref{eq:truss_T}), extended by the scalar rotation DOF.

\begin{figure}[htbp]
\centering
\begin{tikzpicture}[>=Stealth, x=1.1cm, y=1.0cm,
    disp/.style={->,thick,blue!70!black}]

  \def\angledeg{35}

  %% Element
  \draw[line width=2.5pt,gray!55]
        (0,0) -- ({4.5*cos(\angledeg)},{4.5*sin(\angledeg)});
  \filldraw (0,0)                        circle (3.5pt)
        node[below left=1pt] {\small\bfseries 1};
  \filldraw ({4.5*cos(\angledeg)},{4.5*sin(\angledeg)}) circle (3.5pt)
        node[above right=1pt] {\small\bfseries 2};

  %% Global axes
  \draw[->,thick] (0,0) -- ++(2.6,0) node[right] {$x$};
  \draw[->,thick] (0,0) -- ++(0,2.6) node[above] {$y$};

  %% Local axes
  \draw[->,thick,blue!75!black] (0,0)
        -- ++({2.2*cos(\angledeg)},{2.2*sin(\angledeg)}) node[right] {$\bar{x}$};
  \draw[->,thick,blue!75!black] (0,0)
        -- ++({-2.2*sin(\angledeg)},{2.2*cos(\angledeg)}) node[above left] {$\bar{y}$};

  %% Angle arc and label
  \draw (1.2,0) arc[start angle=0, end angle=\angledeg, radius=1.2];
  \node at ({1.5*cos(\angledeg/2)},{1.5*sin(\angledeg/2)}) {$\alpha$};

  %% Global DOFs at node 2
  \draw[->,dashed,thick,red!65!black]
        ({4.5*cos(\angledeg)},{4.5*sin(\angledeg)}) -- ++(1.1,0) node[right] {\small $u_2$};
  \draw[->,dashed,thick,red!65!black]
        ({4.5*cos(\angledeg)},{4.5*sin(\angledeg)}) -- ++(0,1.1) node[above] {\small $v_2$};

\end{tikzpicture}
\caption{Frame element inclined at angle $\alpha$ (CCW from the positive
$x$-axis). Global axes $(x,y)$; local axes $(\bar{x},\bar{y})$ in blue.
$\matr{T}_6$ rotates the global DOFs $(u_i,v_i,\theta_i)$ to local
$(\bar{u}_i,\bar{v}_i,\theta_i)$.}
\label{fig:frame_transform}
\end{figure}

\begin{importantbox}
\textbf{Angle convention:} $\alpha$ is measured counter-clockwise from the
positive $x$-axis to the element axis directed from node~1 to node~2.\\[4pt]
\begin{tabular}{@{}ll@{\quad}l@{}}
Horizontal beam, rightward: & $\alpha=0^\circ$,   & $c=1,\;s=0$\\
Vertical column, upward:    & $\alpha=90^\circ$,  & $c=0,\;s=1$\\
Vertical column, downward:  & $\alpha=-90^\circ$, & $c=0,\;s=-1$
\end{tabular}
\end{importantbox}

The element stiffness and load vector in global coordinates are:
\begin{equation}
    \matr{k}^e = \matr{T}_6\transpose\,\bar{\matr{k}}^e\,\matr{T}_6 ,
    \qquad
    \vect{f}^e = \matr{T}_6\transpose\,\bar{\vect{f}}^e .
    \label{eq:frame_ke_global}
\end{equation}
For vertical columns ($c=0$, $|s|=1$), $\matr{T}_6$ exchanges the axial
and transverse DOF blocks: the column's axial stiffness $EA/L$ appears in
the vertical ($v$-$v$) entries of the global stiffness, and its bending
stiffness $12EI/L^3$ appears in the horizontal ($u$-$u$) entries.

\begin{remark}
Global assembly, boundary condition enforcement, and solution proceed
identically to the truss (\cref{ch:truss}) and beam (\cref{ch:beam})
chapters. The PVW and PMPE derivations for the frame element are direct
extensions of \cref{sec:beam_pvw,sec:beam_pmpe}: replace the $4\times4$
beam integrals with the $6\times6$ frame counterparts after applying
$\matr{T}_6$.
\end{remark}

% =============================================================================
\section{Worked Example: Portal Frame under Combined Loading}
\label{sec:frame_example}
% =============================================================================

\begin{example}[Portal frame --- lateral force and distributed beam load]
\label{ex:frame_portal}

A symmetric portal frame has column height $H = 3000\ \text{mm}$ and beam
span $B = 4000\ \text{mm}$. All members share the same cross-section:
$E = 210{,}000\ \text{N/mm}^2$, $A = 6{,}000\ \text{mm}^2$,
$I = 10{,}000{,}000\ \text{mm}^4$. Both column bases are clamped (fixed).
The frame carries:
\begin{itemize}
    \item a rightward horizontal force $F = 10\ \text{kN}$ at the top of
          the left column (node~2), and
    \item a uniform downward load $q = 5\ \text{N/mm}$ on the beam
          (element~2).
\end{itemize}

\begin{figure}[htbp]
\centering
\begin{tikzpicture}[>=Stealth, x=1.1cm, y=1.0cm,
    force/.style={->,ultra thick,red!75!black},
    disp/.style={->,thick,blue!70!black}]

  %% Frame members (columns: 0--3 in y; beam: 0--4 in x at y=3)
  \draw[line width=2pt,gray!60] (0,0) -- (0,3);   % left column
  \draw[line width=2pt,gray!60] (0,3) -- (4,3);   % beam
  \draw[line width=2pt,gray!60] (4,3) -- (4,0);   % right column

  %% Fixed supports (hatched rectangles + horizontal line)
  \fill[pattern=north east lines,pattern color=gray!70]
        (-0.45,-0.30) rectangle (0.45,0);
  \draw[thick] (-0.45,0) -- (0.45,0);
  \fill[pattern=north east lines,pattern color=gray!70]
        (3.55,-0.30) rectangle (4.45,0);
  \draw[thick] (3.55,0) -- (4.45,0);

  %% Distributed downward load on beam
  \foreach \xp in {0.2,0.5,...,3.8}{
      \draw[disp,->,thin] (\xp,3.65) -- (\xp,3.12);
  }
  \draw[disp,thick] (0,3.65) -- (4,3.65);
  \node[above=2pt,blue!70!black] at (2,3.68) {$q = 5\ \text{N/mm}$};

  %% Lateral force at node 2
  \draw[force] (-1.6,3) -- (-0.08,3)
       node[above=3pt,midway] {$F = 10\ \text{kN}$};

  %% Nodes
  \filldraw (0,0) circle (3pt) node[below left=3pt] {\small\bfseries 1};
  \filldraw (0,3) circle (3pt) node[above left=3pt] {\small\bfseries 2};
  \filldraw (4,3) circle (3pt) node[above right=3pt] {\small\bfseries 3};
  \filldraw (4,0) circle (3pt) node[below right=3pt] {\small\bfseries 4};

  %% Element numbers
  \node[left=10pt]  at (0,1.5) {\small\textcircled{\small 1}};
  \node[above=6pt]  at (2,3)   {\small\textcircled{\small 2}};
  \node[right=10pt] at (4,1.5) {\small\textcircled{\small 3}};

  %% Dimension lines
  \draw[<->,thin] (-0.9,0) -- (-0.9,3)
       node[fill=white,inner sep=1pt,midway,rotate=90]
       {$H = 3000\ \text{mm}$};
  \draw[<->,thin] (0,-0.7) -- (4,-0.7)
       node[fill=white,inner sep=1pt,midway] {$B = 4000\ \text{mm}$};

\end{tikzpicture}
\caption{Portal frame: four nodes, three elements, clamped bases at
nodes~1 and~4. Combined loading: horizontal force $F$ at node~2 and
uniform vertical load $q$ on the beam.}
\label{fig:frame_portal}
\end{figure}

\medskip
\noindent\textbf{DOF ordering:}
$[u_1,v_1,\theta_1,\,u_2,v_2,\theta_2,\,u_3,v_3,\theta_3,\,u_4,v_4,\theta_4]
= \text{DOFs}\ [1\ldots12]$.\\
Clamped DOFs: $\{1,2,3,10,11,12\}=\mathbf{0}$.
Free DOFs: $\{4,5,6,7,8,9\}$.

\medskip
\noindent\textbf{Step 1: Element data.}

\begin{table}[htbp]
\centering
\small
\caption{Element data for the portal frame example.}
\label{tab:frame_elements}
\begin{tabular}{@{}ccccccc@{}}
\toprule
Elem & Nodes & $L$ (mm) & $\alpha$ ($^\circ$) & $c$ & $s$ & Global DOFs \\
\midrule
1 & $1\to2$ & 3000 &  90 & $0$ & $+1$ & 1,2,3,4,5,6 \\
2 & $2\to3$ & 4000 &   0 & $1$ & $0$  & 4,5,6,7,8,9 \\
3 & $3\to4$ & 3000 & $-90$ & $0$ & $-1$ & 7,8,9,10,11,12 \\
\bottomrule
\end{tabular}
\end{table}

\medskip
\noindent\textbf{Step 2: Local stiffness constants.}

Let $a=EA/L$, $b=12EI/L^3$, $d=6EI/L^2$, $e_4=4EI/L$, $e_2=2EI/L$.

\smallskip
For the columns (elements 1 and 3, $L_c = 3000\ \text{mm}$):
\begin{align*}
  a_c &= \frac{210{,}000 \times 6{,}000}{3{,}000} = 420{,}000\ \tfrac{\text{N}}{\text{mm}},
  &b_c &= \frac{12 \times 210{,}000 \times 10^7}{3{,}000^3} = 933.3\ \tfrac{\text{N}}{\text{mm}},\\
  d_c &= \frac{6 \times 210{,}000 \times 10^7}{3{,}000^2} = 1{,}400{,}000\ \text{N},
  &e_{4c} &= \frac{4 \times 210{,}000 \times 10^7}{3{,}000} = 2.800\times10^9\ \text{N\,mm},\\
  &&e_{2c} &= \tfrac{1}{2}\,e_{4c} = 1.400\times10^9\ \text{N\,mm}.
\end{align*}

For the beam (element 2, $L_b = 4000\ \text{mm}$):
\begin{align*}
  a_b &= \frac{210{,}000 \times 6{,}000}{4{,}000} = 315{,}000\ \tfrac{\text{N}}{\text{mm}},
  &b_b &= \frac{12 \times 210{,}000 \times 10^7}{4{,}000^3} = 393.75\ \tfrac{\text{N}}{\text{mm}},\\
  d_b &= \frac{6 \times 210{,}000 \times 10^7}{4{,}000^2} = 787{,}500\ \text{N},
  &e_{4b} &= \frac{4 \times 210{,}000 \times 10^7}{4{,}000} = 2.100\times10^9\ \text{N\,mm},\\
  &&e_{2b} &= 1.050\times10^9\ \text{N\,mm}.
\end{align*}

\medskip
\noindent\textbf{Step 3: Transformation matrices.}

For elements~1 and~3, the column geometry gives $c=0$, $|s|=1$:
\[
    \matr{T}_6^{(1)} =
    \begin{bmatrix}
         0 &  1 & 0 & 0 &  0 & 0 \\
        -1 &  0 & 0 & 0 &  0 & 0 \\
         0 &  0 & 1 & 0 &  0 & 0 \\
         0 &  0 & 0 & 0 &  1 & 0 \\
         0 &  0 & 0 &-1 &  0 & 0 \\
         0 &  0 & 0 & 0 &  0 & 1
    \end{bmatrix}\!,\quad
    \matr{T}_6^{(3)} =
    \begin{bmatrix}
         0 & -1 & 0 &  0 &  0 & 0 \\
         1 &  0 & 0 &  0 &  0 & 0 \\
         0 &  0 & 1 &  0 &  0 & 0 \\
         0 &  0 & 0 &  0 & -1 & 0 \\
         0 &  0 & 0 &  1 &  0 & 0 \\
         0 &  0 & 0 &  0 &  0 & 1
    \end{bmatrix}\!.
\]
For the beam ($\alpha=0^\circ$): $\matr{T}_6^{(2)} = \matr{I}_6$, so the beam's
local and global stiffness matrices are identical.

\end{example}
